\documentclass[12pt]{article}
\usepackage{espd}
\usepackage{amsmath}
\docode{RU.МИФИ.52714-01 СТ 01-1}
% \makeapproval{Наименование}{Вид документа}{Шифр ЛУ}{Шифр}
\newcommand{\univ}{Национальный исследовательский ядерный университет «МИФИ»}
\newcommand{\course}{Дисциплина «Генеративный искусственный интеллект: модели и приложения»}
\newcommand{\progname}{Компьютерная игра «Dreich Day»}
\newcommand{\makeapproval}[3]{%
\thispagestyle{empty}
\begin{center}\singlespacing
\univ\\[2mm]\course
\end{center}
\vspace{10mm}
\noindent\begin{minipage}[t]{0.48\textwidth}\singlespacing
СОГЛАСОВАНО\\[2mm]Студент группы Б25-527\\[6mm]\rule{30mm}{0.4pt} А.\,И.~Лобанов\\[2mm]«\rule{8mm}{0.4pt}»\rule{30mm}{0.4pt} 2026 г.
\end{minipage}\hfill
\begin{minipage}[t]{0.48\textwidth}\singlespacing\raggedleft
УТВЕРЖДАЮ\\[2mm]Преподаватель дисциплины\\[6mm]\rule{30mm}{0.4pt} А.\,А.~Кузьмин\\[2mm]«\rule{8mm}{0.4pt}»\rule{30mm}{0.4pt} 2026 г.
\end{minipage}
\vfill
\begin{center}
{\bfseries\MakeUppercase{\progname}}\\[6mm]
{\bfseries\MakeUppercase{#1}}\\[4mm]
ЛИСТ УТВЕРЖДЕНИЯ\\[4mm]
#2
\end{center}
\vfill
\begin{flushright}\singlespacing
Исполнитель:\\студент группы Б25-527\\Лобанов Алексей Игоревич
\end{flushright}
\vspace{10mm}
\begin{center}Москва 2026\end{center}
\clearpage
\thispagestyle{empty}
\noindent УТВЕРЖДЁН\\#2\\
\vfill
\begin{center}
{\bfseries\MakeUppercase{\progname}}\\[6mm]
{\bfseries\MakeUppercase{#1}}\\[4mm]
#3\\[4mm]
Листов \pageref{LastPage}
\end{center}
\vfill
\begin{center}Москва 2026\end{center}
\clearpage
\setcounter{page}{2}}

\begin{document}
\makeapproval{Спецификация требований}{RU.МИФИ.52714-01 СТ 01-1-ЛУ}{RU.МИФИ.52714-01 СТ 01-1}

\section*{\centering АННОТАЦИЯ}
Настоящий документ содержит спецификацию требований к компьютерной игре «Dreich Day», разработанной в рамках большого домашнего задания по дисциплине «Генеративный искусственный интеллект: модели и приложения».

Спецификация детализирует техническое задание RU.МИФИ.52714-01 ТЗ 01-1: в ней перечислены пронумерованные функциональные и нефункциональные требования, требования к данным, полный перечень скрытых правил жителей и заданий дня, а также критерии приёмки и трассировка требований к итерациям разработки на платформе WebSim. Требования сформулированы проверяемо: для каждого указано, как убедиться в его выполнении.

Документ оформлен в соответствии с ГОСТ 19.101-77 и ГОСТ 19.106-78.
\clearpage
\tableofcontents
\clearpage

\section{Введение}
\subsection{Назначение документа}
Спецификация служит единым перечнем требований, по которому выполнялась генерация кода игры агентом WebSim и проверка каждой версии (плейтест). Номер требования (ФТ-\textit{n}, НФТ-\textit{n}) используется в журнале разработки и в критериях приёмки.

\subsection{Область применения}
Игра — браузерное приложение для одного игрока, жанр — симулятор жизни и головоломка-дедукция с временной петлёй. Опубликована на платформе WebSim, исходный код и документация размещены в публичном репозитории.

\subsection{Термины}
\begin{description}
\item[Петля] — один игровой день с 6:00 до 18:00, после которого день начинается заново.
\item[Условия дня] — случайно выбираемые в 6:00 погода утром и после полудня, знак гадюки Натайр, событие дня и стартовое настроение жителей.
\item[Скрытое правило] — детерминированное правило вида «условие $\rightarrow$ поведение» для конкретного жителя.
\item[Наблюдение] — запись в дневнике о том, что игрок видел срабатывание правила.
\item[Конфликт] — ссора двух персонажей, снижающая их настроение и силу команды.
\item[История] — цепочка шагов, разрешающая конфликт или сюжетную линию; прогресс сохраняется между петлями.
\end{description}

\subsection{Ссылочные документы}
Техническое задание RU.МИФИ.52714-01 ТЗ 01-1; пояснительная записка RU.МИФИ.52714-01 81 01-1; ГОСТ 19.101-77, ГОСТ 19.106-78, ГОСТ 19.201-78, ГОСТ 19.404-79.

\section{Общее описание}
\subsection{Продукт}
Лондонский синоптик Оливер Оверкаст застревает во временной петле в шотландской деревне Глен-Дрих в День святой Бригиты. Игрок разгадывает правила, по которым живут 11 жителей, влияет на них прогнозом в газете и разговорами, разрешает конфликты, разоблачает подкуп судьи и выигрывает финальный матч по шинти против деревни Силлерберн, после чего петля разрывается.

\subsection{Пользователи}
Игроки уютных инди-игр и детективных головоломок, возраст 14+. Специальная подготовка не требуется: правила игры объясняются на экране «Как играть» и подсказками в первой петле.

\subsection{Ограничения и допущения}
\begin{enumerate}
\item Код генерируется агентом платформы WebSim по промптам на естественном языке; объём доработок ограничен суточным лимитом кредитов платформы.
\item Игра работает в браузере без установки и без внешних серверов; прогресс хранится в \texttt{localStorage} браузера.
\item Графика — пиксельные спрайты, нарисованные с помощью ИИ программно (Python, Pillow); звук синтезируется через Web Audio API без звуковых файлов.
\end{enumerate}

\section{Функциональные требования}
\subsection{Цикл дня и временная петля}
\begin{description}
\item[ФТ-1.] Игровой день длится с 6:00 до 18:00; один игровой час равен 50 секундам реального времени. На панели отображаются время, номер петли и значок погоды.
\item[ФТ-2.] В 6:00 каждой петли случайно выбираются условия дня: погода утром и после полудня (солнце, дождь, туман, ветер; после полудня совпадает с утренней с вероятностью 50\,\%), знак Натайр (совпадает с погодой после полудня с вероятностью 70\,\%), событие дня (ярмарка, визит Ангуса, визит Руперта или ничего) и стартовое настроение каждого жителя в диапазоне 40–70.
\item[ФТ-3.] В 18:00 или при выборе «лечь спать» показывается сводка дня, после чего начинается новая петля. Сохраняются только дневник, прогресс историй и конфликтов, выполненные задания и результаты матча.
\end{description}

\subsection{Обряд и прогноз}
\begin{description}
\item[ФТ-4.] В 7:00 жители собираются на площади вокруг норы Натайр, нора остаётся видимой; показывается знак гадюки с пояснением.
\item[ФТ-5.] До 9:00 в редакции газеты игрок выбирает прогноз (солнце, дождь, туман, ветер) один раз за петлю. Жители планируют день по прогнозу.
\item[ФТ-6.] В 12:00 прогноз проверяется: если он сбылся, настроение жителей повышается, если нет — понижается.
\end{description}

\subsection{Жители и скрытые правила}
\begin{description}
\item[ФТ-7.] В деревне 11 жителей. У каждого есть базовое расписание и одно–три скрытых правила, переопределяющих расписание. Правила задаются данными и полностью детерминированы: при одинаковых условиях и прогнозе житель ведёт себя одинаково.
\item[ФТ-8.] Жители перемещаются по общей сетке проходимости с поиском пути и входят в здания только через двери; не стоят на одной клетке.
\item[ФТ-9.] Над жителем отображается значок настроения (пороги 33 и 66). Разговор (клавиша E) показывает реплику по настроению; первый разговор за петлю даёт +5 к настроению.
\end{description}
Перечень реализованных правил приведён в таблице~\ref{tab:rules}.

{\small\singlespacing
\begin{longtable}{|p{12mm}|p{24mm}|p{118mm}|}
\caption{Скрытые правила жителей}\label{tab:rules}\\\hline
Код & Житель & Правило\\\hline\endfirsthead
\hline Код & Житель & Правило\\\hline\endhead
M1 & Мораг & Если утром идёт дождь, Мораг остаётся дома до полудня; пекарня закрыта\\\hline
M2 & Мораг & При прогнозе солнца Мораг выставляет столики у пекарни; дождь в полдень портит ей настроение\\\hline
H1 & Хэмиш & При прогнозе дождя или ветра Хэмиш пропускает рыбалку и с 8:00 сидит в пабе\\\hline
H2 & Хэмиш & В утреннем тумане Хэмиш ничего не ловит и в 10:00 уходит в паб, теряя 15 настроения\\\hline
H3 & Хэмиш & Если после полудня ветрено, Хэмиш уходит в каюту лодки вместо паба\\\hline
F1 & Фергус & Во время ярмарки Фергус держит лоток на площади с 10:00 до 16:00\\\hline
F2 & Фергус & Если Хэмиш приходит в паб до полудня, оба получают +10 к настроению\\\hline
F-R & Фергус & В день визита Руперта настроение Фергуса падает на 15\\\hline
D1 & Дуги & В день приезда Ангуса Дуги запирается в хижине с 10:00 до 15:00\\\hline
D2 & Дуги & Если утром солнечно, тренировка Дуги на поле длится до 13:00\\\hline
D3 & Дуги & После примирения Дуги и Ангус тренируются вместе на поле\\\hline
E1 & Элспет & Если прогноз не сбылся, Элспет с 12:00 до 16:00 пишет поправку в редакции\\\hline
E2 & Элспет & Во время ярмарки Элспет репортирует с площади с 10:00 до 13:00\\\hline
I1 & Иэн & В утреннем тумане Иэн звонит в колокол в 8:00, пятеро жителей приходят к кирку\\\hline
I2 & Иэн & В полдень Иэн навещает самого печального жителя, если его настроение ниже 34\\\hline
R1 & Рэб & Когда Рэб в духе (настроение не ниже 66) и на площади, его волынка раз в час даёт +5 всем в радиусе 5 клеток\\\hline
R2 & Рэб & В ветер Рэб отказывается играть и уходит в паб\\\hline
P1 & Рэб & После примирения с Иэном Рэб в 15:00 играет у дверей кирка; Иэну и Бриди +10\\\hline
K1 & Кенни & При прогнозе дождя Кенни держит стадо на пастбище до 15:00, затем идёт в паб\\\hline
K2 & Кенни & В день ярмарки Кенни приводит двух овец на площадь с 10:00 до 14:00\\\hline
Ma1 & Малкольм & В утреннем тумане горн не тянет: Малкольм с 9:00 до 12:00 в пабе\\\hline
Ma2 & Малкольм & Когда Рэб играет на площади, Малкольм отбивает ритм молотом: +10 раз в день\\\hline
Ro1 & Рори & В солнечный полдень Рори водит туристов вдоль берега и к 17:00 устаёт\\\hline
Ro2 & Рори & При прогнозе дождя Рори отменяет прогулку и отдыхает в сторожке\\\hline
B1 & Бриди & Если Натайр выползла из норы, в 7:30 Бриди объявляет доброе знамение: всем +5\\\hline
B2 & Бриди & В полдень Бриди идёт к жителю с настроением ниже 30: ему +10, ей $-5$\\\hline
\end{longtable}}

\subsection{Дневник}
\begin{description}
\item[ФТ-10.] Дневник открывается клавишей J. Наблюдение записывается автоматически, когда игрок находится рядом с жителем в момент срабатывания правила и житель действительно находится в месте, предписанном правилом.
\item[ФТ-11.] Правило открывается после наблюдения в двух разных петлях; неразгаданные правила отображаются как «???».
\item[ФТ-12.] В дневнике есть разделы «Задания», «Истории», «Финал» (число сыгранных матчей и лучший счёт) и раздел о Руперте «Чужак».
\end{description}

\subsection{Задания дня}
\begin{description}
\item[ФТ-13.] В 6:00 выбирается одно задание из пула, выполнимое при текущих скрытых условиях. Задание объявляется в утренней газете и отображается на панели статуса.
\item[ФТ-14.] Первое выполнение задания открывает подсказку к одному неразгаданному правилу. Прогресс отображается как «Задания: \textit{n}/10».
\end{description}

{\small\singlespacing
\begin{longtable}{|p{24mm}|p{50mm}|p{80mm}|}
\caption{Пул заданий дня}\label{tab:tasks}\\\hline
Заказчик & Задание & Цель\\\hline\endfirsthead
\hline Заказчик & Задание & Цель\\\hline\endhead
Фергус & Скучно в пабе & Хэмиш в пабе до 12:00\\\hline
Мораг & Торговля на свежем воздухе & Столики Мораг на улице в 12:30, после полудня без дождя\\\hline
Хэмиш & Улов к обеду & Хэмиш на причале до 12:00\\\hline
Фергус & Хороший день для деревни & Мораг, Хэмиш и Фергус довольны до 17:00\\\hline
Газета & Честное слово синоптика & Прогноз сбылся и двое довольны в 17:00\\\hline
Дуги & Тренировка & Дуги на поле в 12:00\\\hline
Элспет & Чистый выпуск & Прогноз сбылся, Элспет довольна в 16:00\\\hline
Иэн & Полный кирк & Пятеро жителей у кирка в 9:00\\\hline
Рэб & Волынка на площади & Рэб играет на площади в 11:00\\\hline
Кенни & Овцы на ярмарке & Кенни с овцами на площади в 12:00\\\hline
\end{longtable}}

\subsection{Конфликты и истории}
\begin{description}
\item[ФТ-15.] Конфликты задаются данными: пара персонажей, эффект неразрешённого конфликта, условия и эффект разрешения. Реализованы конфликты «Старые друзья» (Дуги — Ангус), «Пикша и булки» (Мораг — Хэмиш), «Псалмы и волынка» (Иэн — Рэб).
\item[ФТ-16.] История «Старые друзья»: визит Ангуса, разговоры с Ангусом и Дуги, сцена примирения в пабе; результат — совместное правило D3.
\item[ФТ-17.] История «Подкупленный судья»: подслушать разговор Руперта на площади (11:00–13:00), подобрать квитанцию о подкупе у редакции (после 15:00), отдать её Элспет при её настроении не ниже 50. Результат — экстренный выпуск, заголовок «СКУННЕРСБИ ПОДКУПИЛ СУДЬЮ» во всех следующих газетах, Руперт больше не появляется, флаг честного судьи для матча.
\end{description}

\subsection{Финальный матч}
\begin{description}
\item[ФТ-18.] В 15:30 жители с настроением не ниже 40 идут на поле для шинти. Матч начинается в 16:00, если Оливер находится не дальше 4 клеток от поля.
\item[ФТ-19.] Перед матчем показывается лист команды: состав с портретами и настроением, состояние конфликтов, судья, погода и прогноз, полосы силы команд и вердикт.
\item[ФТ-20.] Сила команды рассчитывается по формуле
\[
V = \sum_{i\in\text{вышли}} \mathrm{mood}_i + 60\,n_{\text{разреш}} - 40\,n_{\text{неразреш}} + W + R,\qquad
\rho = \mathrm{clamp}\!\left(\frac{V}{900};\ 0{,}3;\ 1{,}4\right),
\]
где $W$ — поправка на погоду ($+60$ при дожде или тумане после полудня и прогнозе солнца, $-30$ при солнце), $R = -150$, если подкуп судьи не разоблачён. Все числа вынесены в объект конфигурации.
\item[ФТ-21.] Матч — мини-игра 6 на 6 длительностью 75 секунд: Оливер под управлением игрока (WASD — бег, Пробел — пас, E — удар) и пять жителей с наибольшим настроением (Хэмиш — вратарь) против команды Силлерберна с капитаном Ангусом. Скорость, точность и отбор мяча игроков ИИ зависят от $\rho$.
\item[ФТ-22.] Победа засчитывается, если «Гадюки» забили больше и хотя бы один гол забил Оливер. При поражении показывается подсказка о самом слабом модификаторе, затем сводка дня и новая петля.
\end{description}

\subsection{Концовка, меню и обучение}
\begin{description}
\item[ФТ-23.] После победы показывается финальная последовательность из четырёх сцен (свисток, отъезд Руперта, праздник в пабе, утро 2 февраля) и титры с указанием автора, применения генеративного ИИ и числа петель до победы. Факт прохождения сохраняется.
\item[ФТ-24.] Главное меню: «Продолжить» (при наличии сохранения), «Новая игра» (с подтверждением, если есть сохранение), «Как играть».
\item[ФТ-25.] В первой петле показываются подсказки: об обряде (6:05), о прогнозе (7:35), о наблюдении за жителями (9:05), о матче (15:00).
\item[ФТ-26.] Звуки синтезируются через Web Audio: свисток в начале и конце матча, сигналы гола, щелчок кнопок; клавиша M включает и выключает звук.
\end{description}

\section{Нефункциональные требования}
\begin{description}
\item[НФТ-1.] Игра открывается в современном браузере (Chromium, Firefox) без установки и без ошибок в консоли.
\item[НФТ-2.] Анимация плавная; игровой цикл построен на \texttt{requestAnimationFrame}.
\item[НФТ-3.] Весь интерфейс и все реплики — на русском языке; имена и названия — шотландские, «говорящие».
\item[НФТ-4.] Графика — пиксель-арт с единой пасмурной палитрой и тёплым светом как акцентом; спрайты персонажей 16×24 пикселя, масштаб без сглаживания.
\item[НФТ-5.] Прогресс сохраняется в \texttt{localStorage} и восстанавливается при повторном открытии; повреждённое или отсутствующее сохранение не мешает начать новую игру.
\item[НФТ-6.] Правила, задания, конфликты, истории и параметры баланса хранятся как данные, чтобы новые элементы добавлялись без переписывания движка.
\end{description}

\section{Требования к данным}
\begin{description}
\item[Правило:] \texttt{\{id, npc, when(state), effect: \{location, from, to, moodDelta\}, text\}}.
\item[Житель:] имя, должность, лист спрайтов, базовое расписание, место на обряде, реплики по настроению и подсказки к правилам.
\item[Конфликт:] \texttt{\{id, a, b, title, resolved\}} и прогресс шагов истории.
\item[Сохранение:] номер петли, дневник (наблюдения и открытые правила), выполненные задания, прогресс историй, флаг честного судьи, число матчей и лучший счёт, флаг прохождения.
\end{description}

\section{Критерии приёмки и трассировка}
Каждая версия принималась по плейтесту: отсутствие ошибок в консоли и выполнение требований соответствующей итерации. Трассировка требований к итерациям и запросам журнала разработки (приложение к пояснительной записке) приведена в таблице~\ref{tab:trace}.

{\small\singlespacing
\begin{longtable}{|p{30mm}|p{40mm}|p{24mm}|p{56mm}|}
\caption{Трассировка требований}\label{tab:trace}\\\hline
Требования & Итерация & Запросы & Проверка\\\hline\endfirsthead
\hline Требования & Итерация & Запросы & Проверка\\\hline\endhead
ФТ-1 — ФТ-11, НФТ-1 — НФТ-5 & 1. Петля & 1–4 & 6 правил открываются в дневнике, прогноз меняет поведение\\\hline
ФТ-13, ФТ-14 & 2. Цели & 5–6 & задание появляется в газете и засчитывается\\\hline
ФТ-4, ФТ-7 (новые жители), НФТ-4 & 2.2–2.3. Сюжет & 7–12 & 6 жителей на обряде, история «Старые друзья» проходится\\\hline
ФТ-7, ФТ-15 — ФТ-17 & 3. Вся деревня & 13–14 & 11 жителей, все шаги истории «Подкупленный судья»\\\hline
ФТ-18 — ФТ-22 & 4. Матч & 15–17 & идеальная команда побеждает, слабая проигрывает\\\hline
ФТ-23 — ФТ-26, НФТ-5 & 5. Финал и полировка & 18 и далее & концовка и титры достижимы, меню, подсказки, звук\\\hline
\end{longtable}}

\changelist
\end{document}
